\section{Rational primes and Dirichlet characters}
\label{sec:dirichlet}
For integers $x\geq 0$, $q\geq1$ and $0\leq a<q$, set
\[
 M_{q,a}(x)=\sum_{\substack{1\leq n\leq x\\n\bmod q=a}}\mu(n),
 \qquad M(x)=M_{1,0}(x).
\]
Here $\mu(1)=1$. A square prime divisor of $n$ gives $\mu(n)=0$.
Otherwise each distinct prime divisor contributes a factor $-1$.
We call $a$ a reduced residue when $\gcd(a,q)=1$. This deliberately
includes $a=0$ for $q=1$. For example,
$M_{3,1}(10)=\mu(1)+\mu(4)+\mu(7)+\mu(10)=1$,
whereas $M(10)=-1$.

\paragraph{Finite evaluation.}
\texttt{DirichletData.lean} tests primality of $p$ by exhausting
$2\leq d<p$. To evaluate $\mu(n)$ for $n>0$, it starts at $+1$ and
visits $p=1,\ldots,n$. At a prime divisor it replaces the accumulator
by $0$ if $p^2\mid n$ and changes its sign otherwise. After step $p$
the accumulator is the product of the local factors for primes at most
$p$. This invariant follows immediately by the three cases in the
update. Every prime divisor of $n$ is eventually visited. Finally the
recursion for $M_{q,a}(x+1)$ adds $\mu(x+1)$ exactly when
$(x+1)\bmod q=a$. These two finite invariants identify the evaluator
with the displayed sums. The file proves the two recursion branches and
their specialization to $M$. Its convention $\mu(0)=0$ never
contributes. All loops have input-supplied bounds. Signed numbers can
be coded by the parity of a natural, and $\gcd(a,q)$ can equivalently
be computed by retaining the last common divisor up to $q$ for $q>0$.

\paragraph{From zeros to the arithmetic sentences.}
We work externally in classical analytic number theory throughout this
paragraph and the converse below. Put $\theta=1/2+\delta$, where
$\delta>0$. Under RH, Littlewood's estimate is $|M(x)|\ll_\delta
x^\theta$. One can also deduce it from
\href{https://arxiv.org/abs/0705.0723}{Soundararajan, \emph{Partial sums
of the M\"obius function}, Theorem~1}: the extra factor beyond
$\sqrt{x}$ has logarithm $o(\log x)$. This deduction uses RH for
$\zeta$ alone. Under Dirichlet GRH,
\href{https://arxiv.org/abs/1406.7326}{Lynnelle Ye, \emph{Bounding sums
of the M\"obius function over arithmetic progressions}, Theorem~1.1}
gives on reduced residues, for all sufficiently large $x$,
\[
 |M_{q,a}(x)|\ll \sqrt{x}
 \exp\!\left((\log x)^{1/2}(\log\log x)^4\right)
\]
with an absolute implied constant (fix the theorem's parameter at $1$).
The ratio $(\log\log x)^4/\sqrt{\log x}$ tends to zero. Consequently
one integer $B_\delta\geq1$ works simultaneously for every $q,a$ in
$|M_{q,a}(x)|\leq B_\delta x^\theta$ for $x\geq1$. To include the
small values of $x$, use the uniform bound $|M_{q,a}(x)|\leq x$.
This proves the stated uniformity without a search over moduli.
Rounding a classical constant upwards proves existence of an integer
witness in the standard model. It supplies no proof of the conjectured
witnesses in iPRA.

\paragraph{Recovering the zeros.}
In the reverse direction even constants depending on $q,a,\delta$
suffice. Fix a character $\chi$ modulo $q$, extended by zero on
nonunits. Its twisted summatory function is
\[
 A_\chi(x)=\sum_{n\leq x}\mu(n)\chi(n)
 =\sum_{\substack{0\leq a<q\\(a,q)=1}}\chi(a)M_{q,a}(x).
\]
This exact regrouping uses finiteness and periodicity. Conversely the
character orthogonality relation recovers each reduced-residue sum as
$\varphi(q)^{-1}\sum_\chi\overline{\chi(a)}A_\chi(x)$.
In particular the arithmetic assertions control the full complex values,
including characters of order greater than two. The triangle inequality
gives $|A_\chi(x)|\leq B x^\theta$, where one may sum the finitely many
class constants to form $B$. The bound extends to real $x\geq1$ by
taking its integer part.

Choose a compact subset of $\Re s>1/2$. Take $\delta>0$ small enough
that $\theta<\inf\Re s$ on this set. With $\sigma=\Re s$, summation by
parts for $1\leq N<Y$ gives the exact tail
\[
 A_\chi(Y)Y^{-s}-A_\chi(N)N^{-s}
       +s\int_N^Y A_\chi(t)t^{-s-1}\,dt .
\]
Its modulus is at most
$B(2+|s|/(\sigma-\theta))N^{-(\sigma-\theta)}$.
Thus $F_\chi(s)=\sum_{n\geq1}\mu(n)\chi(n)n^{-s}$ is holomorphic
in $\Re s>1/2$ by locally uniform convergence. For $\Re s>1$,
absolute convergence of the Euler product gives
$L(s,\chi)F_\chi(s)=1$. The identity theorem extends this identity
through the half-plane, with $s=1$ removed when $L$ has a pole.
A zero there would make the product vanish. For primitive $\chi$,
the functional equation pairs zeros in $0<\Re s<1$ with zeros of the
conjugate character reflected across $\Re s=1/2$. The gamma factors
are finite and nonzero in this strip. Since the argument applies to
every character, a zero on the left would give a forbidden zero on the
right. This proves Dirichlet GRH. Extra Euler factors for imprimitive
characters can vanish only on $\Re s=0$, outside the open strip.
Taking just the unique character modulo $1$ throughout gives the
converse to the RH estimate, without assuming any other progression
bound.

\paragraph{Precision obtained from a witness.}
For positive integers $h,T$, restrict to $\Re s\geq1/2+1/h$ and
$|s|\leq T$. Request a growth witness at any available slack
$\delta\leq 1/(2h)$ and let $B\geq1$ be the resulting integer
bound for $A_\chi$ (for a uniform progression bound $C$, $B=qC$
suffices). The preceding calculation bounds every tail beyond
\[
 N=\bigl(2^j B(2+2hT)+1\bigr)^{2h}
\]
by $2^{-j}$. This explicit primitive-recursive expression explains what
a supplied arithmetic witness buys at a prescribed precision. The
complex analytic identity theorem and the functional equation in the
correspondence argument are external results.
